%! Inverse Functions — Unit 5, Lesson 5.7 — Warm-Up (Answer Key)
\documentclass[10pt]{article}
\usepackage{algebra2-article}
\usepackage{algebra2-key}   % pulls in -boxes; do NOT also load -boxes
\newcommand{\TallMath}[1]{$\displaystyle #1\rule[-1.4em]{0pt}{3.2em}$}
% In-formula answer slot (\ans is text-mode and must never appear inside $...$).
\newcommand{\mans}[1]{{\color{keyred}\mathbf{#1}}}

\begin{document}

\pageheader{Unit 5, Lesson 5.7}{Warm-Up --- Answer Key}
\namedateperiod

\vspace{0.10in}

\begin{notesbox}{Getting Ready: Skills We'll Use Today}
\small Three short items. Every one of them is a piece of today's lesson --- item 2 is
\emph{yesterday's} discovery, and today it gets its name.

\vspace{4pt}
\textbf{1. Solve for the other variable (Algebra 1 and Unit 5).} Each equation is solved for $y$. Solve it
for $x$ instead.
\begin{center}
\renewcommand{\arraystretch}{2.0}
\begin{tabularx}{0.96\linewidth}{@{}X X X@{}}
(a) \; $y=3x-7$ \quad $x=$ \ans{$\dfrac{y+7}{3}$} &
(b) \; $y=\dfrac{x+5}{2}$ \quad $x=$ \ans{$2y-5$} &
(c) \; $y=x^{3}+1$ \quad $x=$ \ans{$\sqrt[3]{y-1}$} \\
\end{tabularx}
\end{center}
\vspace{-4pt}
Part (c) needed an operation from this unit. Which one, and why was it the right one to use?
\ansline{The cube root --- it is what undoes cubing. Subtract $1$ first, then take
$\sqrt[3]{\phantom{x}}$ of both sides. (No $\pm$ is needed: an odd index has exactly one real root.)}

\vspace{0.12in}

\textbf{2. Both orders (Lesson 5.6).} Let $f(x)=2x+6$ and $g(x)=\dfrac{x-6}{2}$.
\begin{center}
\renewcommand{\arraystretch}{1.8}
\begin{tabularx}{0.96\linewidth}{@{}X X@{}}
(a) \; $(f\circ g)(x)=$ \ans{$2\!\left(\frac{x-6}{2}\right)+6=x-6+6=x$} &
(b) \; $(g\circ f)(x)=$ \ans{$\frac{(2x+6)-6}{2}=\frac{2x}{2}=x$} \\
\end{tabularx}
\end{center}
\vspace{-6pt}
Yesterday you met one or two pairs that did this. In your own words: what are these two functions
doing to each other? \ansline{Each one cancels the other out. Whatever $f$ does to a number, $g$
takes it straight back --- and it works in either order, so the pair returns the $x$ you started
with.}

\vspace{0.12in}

\textbf{3. Undo the number trick.} A classmate says: \emph{``Think of a number. Double it. Then add
$5$. My answer is $17$.''}
\begin{center}
\renewcommand{\arraystretch}{1.6}
\begin{tabularx}{0.96\linewidth}{@{}X c c@{}}
\hline
 & \textbf{First move} & \textbf{Second move} \\
\hline
The trick did this, in this order & double & add $5$ \\
To get back, you do this, in this order & \ans{subtract $5$} & \ans{halve} \\
\hline
\end{tabularx}
\end{center}
\vspace{-2pt}
Their number was \ans{$6$}. \; To undo a two-step process you reverse each step \emph{and} you
also reverse the \ans{order}.

\end{notesbox}

\begin{teachernote}
\small Five minutes. Do not resolve any of the three --- each is picked up by name within twenty
minutes.
\begin{itemize}[leftmargin=1.3em, nosep]
  \item \textbf{Item 1 is the engine of the whole lesson.} ``Swap and solve'' is nothing but
        \emph{solve for the other variable}, and students can already do it. Part (c) is the unit's
        contribution: a cube root is the tool that undoes a cube, which is the first hint that
        \emph{radical} and \emph{power} functions are inverse families --- the loose thread from
        Lesson 5.0.
  \item \textbf{Item 2 is Lesson 5.6's Tier A discovery, promoted to a warm-up.} Collect the wordings
        students invented yesterday (``they cancel,'' ``it goes backwards,'' ``it undoes it'') and
        write two or three on the board. When \textbf{inverse} appears in Section 1, put it beside
        their words --- the definition should look like a formal version of what they already said.
  \item \textbf{Item 3 plants the day's most under-taught idea: reverse the operations \emph{and} the
        order.} Nobody undoes ``double, then add $5$'' by halving first. Point at that answer again
        during the Hook, when half the room writes $\frac59f-32$ for the Fahrenheit formula ---
        the same mistake, in a place where it is much harder to see.
\end{itemize}
\end{teachernote}

\end{document}
